\section{Revise in separate passes}
\label{sec:revision}

A single revision pass rarely works. Argument, evidence, and prose require
different kinds of attention. Revise from the largest decision to the
smallest. Shreyas compares the last prose pass with the final denoising step in
diffusion: it can sharpen a structure that already exists, but it cannot supply
a missing argument.

\subsection{Pass 1: argument}

Read only the material visible to a skimming reader. Start with the title and
abstract, follow the section headings and figure captions, and finish with the
conclusion. Write down the central claim implied by that path. Compare it with
the paper contract from \Cref{sec:start}. If they differ, fix the claim
hierarchy and section order before editing prose.

Check that each section answers one question and that the sequence is
necessary. Remove duplicate motivation and contribution lists. Move secondary
material out of the path of the main result.

Test the argument on three readers. A student should understand the problem and
its importance. An expert in the paper's field should see the technical
distinction and novelty. A reviewer or expert from a neighboring field should
see the broader intellectual consequence. If the framing works for only one of
them, choose an intentionally narrow audience or revise the argument for
broader readers.

\subsection{Pass 2: evidence}

Trace each major claim to the evidence that supports it. Verify every number
and citation against the source. Check assumptions and uncertainty separately.
A mechanism statement needs mechanism evidence; a performance statement needs
a fair comparison.

Record missing controls and unresolved logic plainly. This pass may weaken a
sentence or narrow a contribution. That is a correction, not a stylistic loss.

\subsection{Pass 3: sections and paragraphs}

Ensure each paragraph helps the reader follow one inference. Keep required
setup close to the relevant result. Move implementation detail, complete
proofs, and secondary sweeps to the appendix when the main argument remains
self-contained.

Check transitions by reading the last sentence of one paragraph and the first
sentence of the next. The first should create or answer a real question. Add a
transition only when the logical connection is missing; do not add a generic
recap.

\subsection{Pass 4: sentences and voice}

Make the minimum effective line edit. Keep the writer's vocabulary and
uncertainty. Preserve the cadence when it carries the argument. Leave strong
sentences alone. Untangle a long sentence when it carries several inferences,
but keep it when its clauses form one clear relation.

Cut throat-clearing and inflated claims. Remove qualifiers that add no meaning;
replace abstract importance with a concrete result. Prefer a direct verb to a
weak verb phrase. Repeat the correct technical term instead of rotating
synonyms for variety.

Watch for canned contrasts and rhetorical question-and-answer setups. Repeated
paragraph shapes can sound equally mechanical. These habits flatten distinct
arguments into the same rhythm. Preserve a precise contrast or short sentence
when it carries real scientific or personal force.

\subsection{Pass 5: the rendered paper}

Read the compiled PDF alongside the source. Inspect every page at normal zoom.
Check the figures and tables first, then inspect equations and links. Make sure
headings do not strand a line at the bottom of a page and that floats appear
near the text that interprets them.

Then read the paper aloud or at speaking pace. This catches buried subjects,
repeated syntax, and sentences whose grammar works but whose logic does not.
The final manuscript should sound like a researcher explaining the result to a
sharp colleague.

Stop when another edit would change only the surface. The last pass should
leave the paper easier to verify, not merely more uniform.
